% TODO make hw stylesheet
\documentclass[11pt]{scrartcl}
\usepackage[a4paper, total={6in, 8in}]{geometry}
\usepackage[usenames,dvipsnames,svgnames]{xcolor}
\usepackage[shortlabels]{enumitem}
\usepackage[framemethod=TikZ]{mdframed}
\usepackage{amsmath,amssymb,amsthm}
\usepackage[colorlinks]{hyperref}

\hypersetup{
	colorlinks=true,
	linkcolor=blue,
	filecolor=magenta,      
	urlcolor=magenta,
	pdftitle={MAT 319 HW}
}

\usepackage{microtype}
\usepackage{mathtools}
\usepackage[headsepline]{scrlayer-scrpage}
\usepackage{thmtools}
\usepackage{todonotes}

\addtolength{\textheight}{3.14cm}
\providecommand{\ol}{\overline}
\providecommand{\eps}{\varepsilon}
\providecommand{\half}{\frac{1}{2}}
\providecommand{\dang}{\measuredangle} %% Directed angle
\providecommand{\CC}{\mathbb C}
\providecommand{\FF}{\mathbb F}
\providecommand{\NN}{\mathbb N}
\providecommand{\QQ}{\mathbb Q}
\providecommand{\RR}{\mathbb R}
\providecommand{\ZZ}{\mathbb Z}
\providecommand{\dg}{^\circ}
\providecommand{\ii}{\item}
\providecommand{\alert}[1]{{\sffamily\textbf{\textcolor{blue}{#1}}}}

\reversemarginpar

\mdfdefinestyle{mdbluebox}{roundcorner=10pt,innerbottommargin=9pt,
	linecolor=blue,backgroundcolor=TealBlue!5,}
\declaretheoremstyle[headfont=\sffamily\bfseries\color{MidnightBlue},
mdframed={style=mdbluebox},]{thmbluebox}

\mdfdefinestyle{mdbloobox}{frametitlefont=\bfseries,innerbottommargin=8pt,
	nobreak=true,backgroundcolor=TealBlue!5,linecolor=blue,}
\declaretheoremstyle[headfont=\bfseries\color{MidnightBlue},
mdframed={style=mdbloobox},headpunct={\\[3pt]},postheadspace=0pt,]{thmbloobox}

\mdfdefinestyle{mdredbox}{frametitlefont=\bfseries,innerbottommargin=8pt,
	nobreak=true,backgroundcolor=Salmon!5,linecolor=RawSienna,}
\declaretheoremstyle[headfont=\bfseries\color{RawSienna},
mdframed={style=mdredbox},headpunct={\\[3pt]},postheadspace=0pt,]{thmredbox}

\mdfdefinestyle{mdgreenbox}{linecolor=ForestGreen,backgroundcolor=ForestGreen!5,
	linewidth=2pt,rightline=false,leftline=true,topline=false,bottomline=false,}
\declaretheoremstyle[headfont=\bfseries\sffamily\color{ForestGreen!70!black},
mdframed={style=mdgreenbox},headpunct={ --- },]{thmgreenbox}

\mdfdefinestyle{mdorangebox}{linecolor=RedOrange,backgroundcolor=RedOrange!5,
	linewidth=2pt,rightline=false,leftline=true,topline=false,bottomline=false,}
\declaretheoremstyle[headfont=\bfseries\sffamily\color{RedOrange!70!black},
mdframed={style=mdorangebox},headpunct={ --- },]{thmorangebox}

\mdfdefinestyle{mdblackbox}{linecolor=black,backgroundcolor=RedViolet!5!gray!5,
	linewidth=3pt,rightline=false,leftline=true,topline=false,bottomline=false,}
\declaretheoremstyle[mdframed={style=mdblackbox}]{thmblackbox}

\declaretheoremstyle[spaceabove=6pt,spacebelow=6pt]{basehead}

\declaretheorem[style=basehead,name=Problem]{problem}
\declaretheorem[style=basehead,name=Problem,numbered=no]{problem*}
\declaretheorem[style=thmbloobox,name=Setup]{setup} % for config geo
\declaretheorem[style=thmredbox,name=Example,sibling=problem]{example}
\declaretheorem[style=thmredbox,name=$\bigstar$ Problem,sibling=problem]{niceprob}
\declaretheorem[style=thmbluebox,name=Theorem,sibling=problem]{theorem}
\declaretheorem[style=thmbluebox,name=Corollary,sibling=theorem]{corollary}
\declaretheorem[style=thmbluebox,name=Proposition,sibling=theorem]{prop}
\declaretheorem[style=thmbluebox,name=Lemma,numberwithin=problem]{lemma}
\declaretheorem[style=thmbluebox,name=Theorem,numbered=no]{theorem*}
\declaretheorem[style=thmbluebox,name=Lemma,numbered=no]{lemma*}
\declaretheorem[style=thmgreenbox,name=Claim,numberwithin=problem]{claim}
\declaretheorem[style=thmgreenbox,name=Claim,numbered=no]{claim*}
\declaretheorem[style=thmblackbox,name=Remark,numberwithin=problem]{remark}
\declaretheorem[style=thmblackbox,name=Remark,numbered=no]{remark*}
\declaretheorem[style=thmgreenbox,name=Definition,sibling=theorem]{definition}
\declaretheorem[style=thmgreenbox,name=Definition,numbered=no]{definition*}
\declaretheorem[style=thmorangebox,name=Warning,sibling=theorem]{warning}
\declaretheorem[style=thmorangebox,name=Warning,numbered=no]{warning*}
\declaretheorem[style=thmorangebox,name=Note,numbered=no]{note}
\declaretheorem[style=thmblackbox,name=Exercise,sibling=problem]{exercise}
\declaretheorem[style=thmblackbox,name=Exercise,numbered=no]{exercise*}
\declaretheorem[style=thmblackbox,name=Question,sibling=problem]{question}
\declaretheorem[style=thmblackbox,name=Question,numbered=no]{question*}

\newenvironment{soln}{\begin{proof}[Solution]}{\end{proof}}
\newenvironment{walkthrough}{\noindent \textbf{Walkthrough.}}{\medskip}
\newlist{walk}{enumerate}{3}
\setlist[walk]{label=\bfseries (\alph*)}

\providecommand{\pow}{\operatorname{Pow}}
\providecommand{\vocab}[1]{{\textbf{\textcolor{ForestGreen}{#1}}}}
\providecommand{\lcm}{\operatorname{lcm}}

\usepackage{asymptote}
\begin{asydef}
	defaultpen(fontsize(10pt));
	unitsize(1cm);
	size(8cm); // set a reasonable default
	usepackage("amsmath");
	usepackage("amssymb");
	settings.tex="pdflatex";
	settings.outformat="pdf";
	// Replacement for olympiad+cse5 which is not standard
	import geometry;
	// recalibrate fill and filldraw for conics
	void filldraw(picture pic = currentpicture, conic g, pen fillpen=defaultpen, pen drawpen=defaultpen)
	{ filldraw(pic, (path) g, fillpen, drawpen); }
	void fill(picture pic = currentpicture, conic g, pen p=defaultpen)
	{ filldraw(pic, (path) g, p); }
	// some geometry
	pair foot(pair P, pair A, pair B) { return foot(triangle(A,B,P).VC); }
	pair orthocenter(pair A, pair B, pair C) { return orthocentercenter(A,B,C); }
	pair centroid(pair A, pair B, pair C) { return (A+B+C)/3; }
	// cse5 abbreviations
	path CP(pair P, pair A) { return circle(P, abs(A-P)); }
	path CR(pair P, real r) { return circle(P, r); }
	pair IP(path p, path q) { return intersectionpoints(p,q)[0]; }
	pair OP(path p, path q) { return intersectionpoints(p,q)[1]; }
	path Line(pair A, pair B, real a=0.6, real b=a) { return (a*(A-B)+A)--(b*(B-A)+B); }
	// cse5 more useful functions
	picture CC() {
		picture p=rotate(0)*currentpicture;
		currentpicture.erase();
		return p;
	}
	pair MP(Label s, pair A, pair B = plain.S, pen p = defaultpen) {
		Label L = s;
		L.s = "$"+s.s+"$";
		label(L, A, B, p);
		return A;
	}
	pair Drawing(Label s = "", pair A, pair B = plain.S, pen p = defaultpen) {
		dot(MP(s, A, B, p), p);
		return A;
	}
	path Drawing(path g, pen p = defaultpen, arrowbar ar = None) {
		draw(g, p, ar);
		return g;
	}
\end{asydef}

\usepackage{mathrsfs}

\ihead{\footnotesize MAT 319 HW01}
\ohead{\footnotesize \today}

\title{\vspace{-2em}MAT 319 HW01}
\author{\large Aditya Pahuja}
\date{\large Due 9/4}
\begin{document}
\maketitle
\begin{problem*}[1.5]
	Prove $1 + \half + \frac 14 + \cdots + \frac{1}{2^n} = 2 - \frac1{2^n}$
	for all positive integers $n$.
\end{problem*}
\begin{soln}
	We show this by induction on $n$.
	For the base case of $n = 1$, we get the obviously true equation
	\[\frac 32 = 1 + \frac{1}{2^1} = 2 - \frac{1}{2^1} = \frac{3}{2}.\]
	
	Assume as the inductive hypothesis that
	\[1 + \half + \frac 14 + \cdots + \frac{1}{2^n} = 2 - \frac1{2^n}\]
	for some $n\in\NN$. Then, we want to show that the equation holds
	when $n$ is replaced with $n + 1$, i.e.
	\[1 + \half + \frac 14 + \cdots + \frac{1}{2^{n + 1}} =
	2 - \frac1{2^{n + 1}}.\]
	This can be done by manipulating the left-hand side as follows:
	\begin{align*}
		\left(1 + \half + \frac 14 + \cdots + \frac{1}{2^n}\right)
		+ \frac{1}{2^{n+1}} &= 2 - \frac1{2^n} + \frac{1}{2^{n+1}}\\
		&= 2 - \frac{2}{2^{n+1}} + \frac{1}{2^{n+1}}\\
		&= 2 - \frac{1}{2^{n+1}}.
	\end{align*}
	
	Therefore, by the principle of mathematical induction,
	\[1 + \frac12 + \frac{1}{4} + \cdots + \frac{1}{2^n} = 2 - \frac{1}{2^n}\]
	for all positive integers $n$.
\end{soln}

\begin{problem*}[2.3]
	Show $\sqrt{2+\sqrt2}$ is not a rational number.
\end{problem*}
We present two solutions.
\begin{proof}[Solution 1]
	Suppose for the sake of contradiction $\sqrt{2 + \sqrt2}$ is rational.
	Then, it can be expressed as a ratio of integers $\frac ab$ with $b\ne 0$.
	By squaring and subtracting two, we find that
	\[\sqrt 2 = \frac{a^2}{b^2} - 2 = \frac{a^2 - 2b^2}{b^2}\]
	which is rational because $a$ and $b$ being integers implies that
	$a^2 - 2b^2$ and $b^2$ are both integers.
	However, this is a contradiction because $\sqrt 2$ is irrational,
	so $\sqrt{2 + \sqrt 2}$ must be irrational.
\end{proof}
\begin{proof}[Solution 2]
	Observe that $\sqrt{2 + \sqrt2}$ is a solution to the polynomial equation
	\[(x^2 - 2)^2 - 2 = 0.\]
	Expanding the left-hand side gives
	$x^4 - 4x^2 + 4 - 2 = x^4 - 4x^2 + 2 = 0$.
	By the rational root theorem (Corollary 2.3 in the textbook),
	any rational value of $x$ satisfying the equation
	must be an integer dividing the constant term $-2$.
	Thus, if $\sqrt{2 + \sqrt2}$ is rational, then, being a root
	of $x^4 - 4x^2 + 2$, it must be an integer.
	However, $1 < \sqrt{2 + \sqrt2} < 2$ because
	\[1 < 2 + \sqrt 2 < 4 \iff -1 < \sqrt 2 < 2,\]
	and the latter inequality is true,
	so $\sqrt{2 + \sqrt2}$ is bounded between two consecutive integers and thus
	not an integer, hence it is irrational.
\end{proof}

\begin{problem*}[3.4]
	Prove (v) and (vii) of Theorem 3.2.
\end{problem*}
\begin{soln}
	We can prove (v) just by taking $a = 1$ in property (iv),
	as $1 = 1^2 \ge 0$. We do have to check that $1^2 \ne 0$;
	this follows from property (vi) of Theorem 3.1,
	as $1\cdot 1 = 0$ would imply $1 = 0$, but obviously this is not true.
	
	To prove (vii), we notice that $a$ and $b$ being positive means
	$\frac 1a$, $\frac 1b$, and $\frac 1{ab}$ are all positive too
	due to properties (iii) and (vi) in 3.2. This means that
	\[0\cdot \frac 1{ab} < a\cdot \frac 1{ab} < b\cdot \frac{1}{ab}
	\iff 0 < \frac 1a < \frac 1b,\]
	as desired.
\end{soln}

\begin{problem*}[3.6]\phantom{a}
	\begin{enumerate}[\textbf{(\alph*)}]
		\ii Prove $|a + b + c| \le |a| + |b| + |c|$
		for all $a$, $b$, $c\in \RR$.
		\ii Use induction to prove
		\[|a_1 + a_2 + \cdots + a_n| \le |a_1| + |a_2| + \cdots + |a_n|\]
		for $n$ real numbers $a_1$, $a_2$, \dots, $a_n$.
	\end{enumerate}
\end{problem*}
\begin{soln}
	For part (a), apply triangle inequality twice:
	\[(|a| + |b|) + |c| \ge |a + b| + |c| \ge |a + b + c|.\]
	
	For part (b), we induct on $n$, as instructed,
	with the base case of $n = 1$ being the trivial statement $|a_1|\le |a_1|$.
	
	Now, assume as the inductive hypothesis that
	\[|a_1 + a_2 + \cdots + a_n| \le |a_1| + |a_2| + \cdots + |a_n|.\]
	We wish to prove that
	\[|a_1 + a_2 + \cdots + a_{n+1}| \le |a_1| + |a_2| + \cdots + |a_{n+1}|,\]
	which is done exactly as in part (a):
	\[|a_1| + |a_2| + \cdots + |a_{n}| + |a_{n+1}| \ge
	|a_1 + a_2 + \cdots + a_n| + |a_{n+1}| \ge |a_1 + a_2 + \cdots + a_{n+1}|,\]
	with the first inequality following from the hypothesis
	and the second from the two-variable version of the inequality.
	
	We have hence proven that $|a_1 + \cdots + a_n| \le |a_1| + \cdots + |a_n|$
	for all positive integers $n$ using the principle of mathematical induction.
\end{soln}

\begin{remark*}
	The triangle inequality can also be proven non-inductively
	by squaring both sides and subtracting $a_1^2 + a_2^2 + \cdots + a_n^2$
	to get
	\[|a_1 + a_2 + \cdots + a_n| \le |a_1| + |a_2| + \cdots + |a_n|
	\iff \sum_{1 \le i < j \le n} 2a_ia_j \le \sum_{1\le i<j \le n} 2|a_ia_j|\]
	which is obviously true because $|x| \ge x$ for all real $x$.
\end{remark*}

\begin{problem*}[4.11]
	Consider $a$, $b\in \RR$ where $a < b$. Use the denseness of $\QQ$ in $\RR$
	to show there are infinitely many rationals between $a$ and $b$.
\end{problem*}
\begin{soln}
	The main idea is to split $(a, b)$ into infinitely many
	disjoint open intervals, each of which contains a rational number
	by denseness.
	The rest of this solution can be seen as a verification
	of the fact that this splitting is possible.
	
	Define the sequence of real numbers $(r_k)$ by $r_0 = a$ and
	\[r_{k+1} = \frac{r_k + b}{2}\]
	for integers $k\ge 0$.
	
	\begin{claim*}
		For all $k$, $r_k < r_{k+1}$.
		Additionally, the set $S = \{r_0, r_1, r_2, r_3, \dots\}$ is
		bounded below by $a$ and bounded above by $b$.
	\end{claim*}
	\begin{proof}
		We prove by induction on $k$ that $r_k < r_{k+1} < b$ for all
		nonnegative integers $k$.
		For the base case of $k = 0$, we already know that $r_0 = a < b$,
		and then $r_0 < r_1 < b$ translates to
		\[a < \frac{a + b}{2} < b \iff 2a = (a + a) < a + b < b + b \]
		We then want to show that if $r_k < r_{k+1} < b$,
		then $r_{k+1} < r_{k+2} < b$.
		Using the definition of $r_k$, the inequality that we want to prove is
		\[\frac{r_k + b}{2} < \frac{r_{k+1} + b}{2} < b,\]
		which is equivalent to the inductive hypothesis after multiplying by $2$
		and subtracting $b$.
		Thus, by the principle of mathematical induction,
		$r_k < r_{k+1} < b$ for all $k$.
		This proves that $b$ is an upper bound for $S$
		and that $(r_k)$ is an increasing sequence.
		
		Finally, every element of $S$ is at least as large as $r_0 = a$
		due to the previous paragraph, so $a = \min S$
		is a lower bound for $S$.
	\end{proof}
	
	Now, we are done, since, by density, there is a rational number
	in $(r_k, r_{k+1})\subset(a, b)$ for every nonnegative integer $k$
	and the $(r_k, r_{k+1})$ intervals are pairwise disjoint
	(since the sequence is increasing)
	which gives us infinitely many rational numbers in the interval $(a, b)$
	as desired.
\end{soln}

\begin{problem*}[4.12]
	Let $\mathbb I$ be the set of irrational numbers. Prove that
	if $a < b$, then there exists $x\in\mathbb I$ such that $a < x < b$.
\end{problem*}
\begin{soln}
	Let $d$ be a rational number satisfying $0 < d < b - a$,
	which must exist by denseness.
	By the archimedean property, we can find an integer $n$ such that
	$a - \sqrt 2 < nd$. There must be a smallest such integer $n$
	because the right-hand side is not bounded below as $n$ decreases;
	let $n_0$ be this smallest value
	(in particular the inequality fails when $n \le \frac{a - \sqrt 2}{d}$,
	so $n_0$ is the smallest integer bigger than $\frac{a - \sqrt2}{2}$).
	Then, we claim that $n_0d < b - \sqrt 2$.
	We prove this by contradiction: if the inequality were false,
	then we find that
	\[(n_0 - 1)d < a - \sqrt2 < b - \sqrt2 \le n_0d,\]
	which implies the false statement
	\[d = n_0d - (n_0 - 1)d > (b - \sqrt2) - (a - \sqrt2) > b - a.\]
	
	Therefore, we have found an integer $n_0$ such that
	\[a < \sqrt2 + n_0d < b.\]
	Since $\sqrt2$ is irrational, so is $\sqrt2 + n_0d$
	(because $n_0d$ is rational), so we have found an irrational number
	between $a$ and $b$.
\end{soln}

\begin{problem*}[5.1]
	Write the following sets in interval notation:
	\begin{enumerate}[\textbf{(\alph*)}]
		\ii $\{x\in\RR : x < 0\}$
		\ii $\{x\in\RR : x^3 \le 8\}$
		\ii $\{x^2 : x\in\RR\}$
		\ii $\{x\in\RR : x^2 < 8\}$
	\end{enumerate}
\end{problem*}
\begin{soln}
	\textbf{(a)} is the negative reals $(-\infty, 0)$.
	
	\textbf{(b)} is $(-\infty, 2]$ because the inequality is equivalent to
	$x \le 2$.
	
	\textbf{(c)} is the nonnegative reals, since the square of a real number
	is always at least zero and all nonnegative real numbers are squares,
	so the corresponding interval is $[0, \infty)$.
	
	\textbf{(d)} is $(-\sqrt8, \sqrt8)$, since the inequality is equivalent to
	$|x| < \sqrt 8$.
\end{soln}

\begin{problem*}[5.2]
	Give the infimum and supremum of each set listed in Exercise 5.1.
\end{problem*}
\begin{soln}
	Reading off the intervals:
	\begin{enumerate}[\textbf{(\alph*)}]
		\ii infimum is $-\infty$, supremum is $0$
		\ii infimum is $-\infty$, supremum is $2$
		\ii infimum is $0$, supremum is $\infty$
		\ii infimum is $-\sqrt8$, supremum is $\sqrt8$
	\end{enumerate}
\end{soln}
\end{document}